\usepackage{amsmath, amssymb, amsfonts, amsthm}
\allowdisplaybreaks
\usepackage[mode=buildnew]{standalone}

\usepackage{geometry}
\geometry{a4paper, margin=1in}
% \usepackage{mdframed}
\usepackage{enumitem}
\setlist{nosep}
\usepackage{booktabs}
\usepackage[style=alphabetic, backend=biber, sorting=ynt]{biblatex}
\usepackage{longtable,array,booktabs}
\usepackage[table]{xcolor}

\usepackage{multicol}
\usepackage{subcaption}
\usepackage{minted}
\setminted{
    baselinestretch=1.0,       % 行间距
    frame=single,               % 边框样式
    framesep=2mm,              % 边框间距
    rulecolor=\color{black},   % 边框颜色
    breaklines,           % 自动换行
    breakanywhere=false,       
    showspaces=false,          % 不显示空格
    showtabs=false,            % 不显示制表符
    tabsize=4,                 % 制表符大小
    autogobble=true,
}
\setmintedinline{bgcolor=yellow!20}

\usepackage{graphicx}
\graphicspath{{images/}}
% \usepackage{subfigure}
\usepackage[most]{tcolorbox} 
\usepackage{tikz}
\usepackage{pgfplots}
\pgfplotsset{compat=1.18}   % 使用较新版本语法
\usetikzlibrary{positioning, arrows.meta, intersections, calc, fit, backgrounds}
\usepgfplotslibrary{fillbetween}  % 用于填充区域

\usepackage{hyperref}
\usepackage{nameref}
\usepackage{cleveref}

\title{\Huge \bfseries 人工智能基础知识手册}
\author{臧炫懿\and DeepSeek-V4-Pro}
\date{\today}

\addbibresource{references.bib}

\renewcommand{\bibfont}{\footnotesize}

\tcbset{
  theorem-style/.style={
    enhanced,
    breakable,                 % 允许跨页
    boxrule=0.4pt,             % 细边框，更素净
    arc=3pt,                   % 圆角
    left=6pt, right=6pt,
    top=4pt, bottom=4pt,
    fonttitle=\bfseries,       % 标题粗体
    attach title to upper,     % 标题与内容在同一框中
    after title={.\ },         % 标题后跟句点
  },
  note-style/.style={
    enhanced,
    breakable,
    boxrule=0.4pt,
    arc=3pt,
    left=6pt, right=6pt,
    top=4pt, bottom=4pt,
    fonttitle=\bfseries,
  }
}

\newtcbtheorem[number within=section]{theorem}{定理}{%
  theorem-style,
  colback=blue!5,              % 极淡蓝背景
  colframe=blue!50!black,      % 深蓝边框
  coltitle=blue!50!black,      % 深蓝标题文字
}{thm}

\newtcbtheorem[use counter from=theorem]{lemma}{引理}{%
  theorem-style,
  colback=green!5,
  colframe=green!50!black,
  coltitle=green!50!black,
}{lm}

\newtcbtheorem[number within=section]{proposition}{命题}{%
  theorem-style,
  colback=purple!5,
  colframe=purple!50!black,
  coltitle=purple!50!black,
}{prop}

\newtcbtheorem[use counter from=theorem]{corollary}{推论}{%
  theorem-style,
  colback=orange!5,
  colframe=orange!50!black,
  coltitle=orange!50!black,
}{cor}

\newtcbtheorem[number within=section]{definition}{定义}{%
  theorem-style,
  colback=red!5,
  colframe=red!50!black,
  coltitle=red!50!black,
}{def}

\newtcbtheorem[number within=section]{example}{例}{%
  theorem-style,
  colback=cyan!5,
  colframe=cyan!50!black,
  coltitle=cyan!50!black,
}{ex}

\newtcolorbox{intro}{
  note-style,
  colback=green!5,
  colframe=green!50!black,
  coltitle=white,
  title={本章当中，你将学到},
}

\tcolorboxenvironment{proof}{
blanker,breakable,left=5mm,
before skip=10pt,after skip=10pt,
borderline west={1mm}{0pt}{red}}


\newtcolorbox{framed}[1]{
  note-style,
  colback=yellow!5,             
  colframe=yellow!50!black,
  coltitle=white,
  title={#1}
}

\tcbset{myformula/.style={colback=yellow!10!white,colframe=red!50!black,
every box/.style={highlight math style={colback=LightBlue!50!white,colframe=Navy}}
}}